\documentclass{article}

\usepackage{xcolor}
\usepackage[colorlinks=true,linkcolor=blue!70!black]{hyperref}
\usepackage{amsmath}
\usepackage{amssymb}
\usepackage{lpalign}

% Global style setters exercised here on purpose.
\setvarsstyle{\scriptscriptstyle}
\setnamedsubeqstyle{\itshape\roman}

\begin{document}

\section*{lpalign Regression}

This file is a compact smoke test for the most fragile \textsf{lpalign} paths:
manual tags, subtags, default subtags, objective-row labels, box forms,
\verb|\downtag|, and a nested matrix/aligned case.

\subsection*{Manual Tags In Ordinary Rows}

\begin{lpregion}
  x &\ge 0 \lptag{R1}\label{reg:R1} \\
  y &\ge 0 \label{reg:R2}
\end{lpregion}

Row \ref{reg:R1} is manually tagged; the next row should remain \ref{reg:R2}
with no skipped counter.

\subsection*{Objective-Row Subtag}

\begin{namedsubeqs}{NLP}[style=\itshape\roman]
  \begin{lpalign}{Minimize}
    {f(\mathbf{x}) \lpsubtag{\textsc{obj}}\label{nlp:obj}}
    g_i(\mathbf{x}) &\ge 0 & \forall i \in I \label{nlp:1} \\
    h_j(\mathbf{x}) &= 0   & \forall j \in J \label{nlp:2} \\
    \mathbf{x}      &\ge 0                    \label{nlp:3}
  \end{lpalign}
\end{namedsubeqs}

Objective row: \eqref{nlp:obj}. Constraint range:
(\ref{nlp:1}--\lpsubref{nlp:3}). Nonnegativity:
\lpsubeqref{nlp:3}. Objective subref: \lpsubeqref{nlp:obj}.

\subsection*{Default Subtags}

\begin{namedsubeqs}{DEF}
\begin{lpregion}
  a &= b \defaultsubtag \label{def:a} \\
  c &= d \defaultsubtag \label{def:b}
\end{lpregion}
\end{namedsubeqs}

These should advance normally: \eqref{def:a}, \eqref{def:b},
with subrefs \lpsubref{def:a} and \lpsubref{def:b}.

\subsection*{Mixed Tag Overrides}

\begin{namedsubeqs}{MIX}
\begin{lpregion}
  p &= q \lptag{M}\lpsubtag{1}\label{mix:1} \\
  r &= s \lptag{S}\defaultsubtag\label{mix:2}
\end{lpregion}
\end{namedsubeqs}

Mixed rows should read \eqref{mix:1} and \eqref{mix:2}, with subrefs
\lpsubref{mix:1} and \lpsubref{mix:2}.

\subsection*{Traditional subequations}

\begin{subequations}\label{subeq:block}
\begin{lpregion}
  u &= v \lpsubtag{\texttt{a}}\label{subeq:a} \\
  w &= z \defaultsubtag\label{subeq:b}
\end{lpregion}
\end{subequations}

Traditional subequations should still work: \eqref{subeq:a},
\eqref{subeq:b}, and subrefs \lpsubref{subeq:a}, \lpsubref{subeq:b}.

\subsection*{Box Forms}

\[
\begin{lpregionbox}[flush=l]
  \mathcal{P} = \left\{
  \begin{lpregionbox}[flush=l]
    A\mathbf{x} &\le \mathbf{b} \\
    \mathbf{x}  &\ge 0
  \end{lpregionbox}
  \right.
\end{lpregionbox}
\]

This section only checks that nested box forms typeset cleanly.

\subsection*{Nested Matrix / Aligned}

\begin{namedsubeqs}{MAT}
\begin{lpalign}[flush=a]{Minimize}
  {\begin{aligned}
      c^\top x &+ d^\top y \\
               &+ e^\top z
    \end{aligned}
    \defaultsubtag\label{mat:obj}}
  \begin{bmatrix}
    1 & 0 \\
    0 & 1
  \end{bmatrix}
  \mathbf{x}
  &=
  \mathbf{b}
  \label{mat:1}
\end{lpalign}
\end{namedsubeqs}

Nested structures should not disturb numbering: \eqref{mat:obj}, \eqref{mat:1}.

\subsection*{Long Row With \texttt{\textbackslash downtag}}

\begin{lpalign}[row-sep=-0.2em]{Maximize}
  {c_1x_1+c_2x_2+c_3x_3+c_4x_4+c_5x_5+c_6x_6+c_7x_7+c_8x_8+c_9x_9+c_{10}
   \span\downtag}
  a_1x_1+a_2x_2+a_3x_3+a_4x_4+a_5x_5+a_6x_6+a_7x_7+a_8x_8
  &\ge \mathbf{b} \downtag \\
  x &\ge 0 \label{down:1}
\end{lpalign}

This is only an escape-hatch test. The numbered short row should remain
referenceable as \eqref{down:1}.

\end{document}
